%% Section 7: N\'{e}ron-Tate distance between points on curves

The goal of this section is to prove Proposition~\ref{PropAlgPtFar}, below.
Namely we will show that $\IQbar$-points on smooth curves are
rather ``sparse'', in the sense that the N\'{e}ron--Tate distance
between two $\IQbar$-points on a smooth curve $C$ is in
general large compared with the Weil height of $C$. %% Now we
%% shall prove the following equivalently form of
%% Proposition~\ref{PropNTDistance}.
%\Gcomment{I think it's better not to assume $S$ to be open, as we use induction on $\dim S$ to prove the proposition. Maybe good to relax the assumption on $S$ in Section 3.}


We use the notation from $\mathsection$\ref{subsec:univfamily}. Recall
that we have fixed a projective compactification
$\overline{\mathbb{A}_g}$ of $\mathbb{A}_g$ over $\IQbar$, an ample
line bundle $\cM$, and a height function
$h_{\overline{\mathbb{A}_g},\cM} \colon
\mathbb{A}_g(\IQbar)\rightarrow \IR$ attached to this pair.
We also fixed a closed immersion $\mathfrak{A}_g$ into
$\IP^n_\IQbar\times \mathbb{A}_g$ over $\mathbb{A}_g$ and
let $\tau\colon \mathbb{M}_g\rightarrow\mathbb{A}_g$ denote the
Torelli morphism. 
If $s\in \mathbb{M}_g(\IQbar)$ then $\mathfrak{C}_s$
is a smooth curve of genus $g$ defined over $\IQbar$.
Moreover, if $P,Q\in \mathfrak{C}_s(\IQbar)$, then $P-Q$ is a
well-defined element of $\mathfrak{A}_g(\IQbar)$ and so is its
N\'eron--Tate height  $\hat h(P-Q)$, see (\ref{eq:NTfinemoduli}).

\begin{proposition}\label{PropAlgPtFar}
  Let $S$ be an irreducible closed subvariety of $\mathbb{M}_g$ defined over $\IQbar$. There
  exist positive constants $c_1, c_2, c_3, c_4$ depending on the choices
  made above and on $S$ with the following property. Let $s\in
  S(\IQbar)$ with $h_{\overline{\mathbb{A}_g},\cM}(\tau(s))
  \ge c_1$. There exists a subset $\Xi_s \subset \mathfrak{C}_s(\IQbar)$
  with $\#\Xi_s\le c_2$ such that any $P \in \mathfrak{C}_s(\IQbar)$
  satisfies the following alternative.
  \begin{enumerate}
  \item[(i)] Either $P \in \Xi_s$;
  \item[(ii)] or $\#\{Q \in \mathfrak{C}_s(\IQbar) : \hat{h}(Q-P) \le h_{\overline{\mathbb{A}_g},\cM}(\tau(s))/c_3\} < c_4$.
  \end{enumerate}
\end{proposition}

%% Let us explain why Proposition~\ref{PropNTDistance} holds, assuming
%% the truth of Proposition~\ref{PropAlgPtFar}. By introducing level
%% structure, we can identify $C$, as in
%% Proposition~\ref{PropNTDistance}, with the fiber of $\mathbb{M}_g$
%% above some $s \in\mathbb{M}_{g}(\IQbar)$. By forgetting the level
%% structure we obtain a morphism
%% $\mathbb{A}_{g}\rightarrow \mathbb{A}_{g,1}$ taking $\tau(s)$ to
%% $[\jac(C)]$. To recover Proposition~\ref{PropNTDistance} it suffices
%% to prove a height lower bound $h(\iota([\jac(C)])) \le c
%% h_{\overline{\mathbb{A}_g},\cM}(\tau(s)) + c'$, where $c>0,c'>0$ are
%% fixed. Both heights arise using immersions of $\mathbb{A}_{g,1}$ and
%% $\mathbb{A}_{g}$ into some projective space. The desired height
%% inequality follows from basic height properties that are essentially
%% consequences of the triangle inequality.

%by the Height Machine.%!!!TODO: ADAPT 1.4. 
%% \Pcomment{TODO: say something about the different heights in \ref{PropAlgPtFar} and \ref{PropNTDistance}}. \Gcomment{In the current version, it's enough to say this, right?}

\begin{proof}
  We fix an immersion of $\mathbb{M}_g$ into
  some projective space and let $\overline{\mathbb{M}_g}$ denote its
  Zariski closure. 
% Fix a compactification $\overline{\mathbb{M}_g}$ of $\mathbb{M}_g$
% and an ample line bundle $\mathcal{N}$ on $\overline{\mathbb{M}_g}$.
% The Height Machine provides a height function
% $h_{\overline{\mathbb{M}_g},\mathcal{N}} \colon
% \overline{\mathbb{M}_g}(\IQbar) \rightarrow \IR$.  
By a standard triangle inequality estimate, there exist constants
$c''>0$ and $c'''\ge 0$ such that
\begin{equation}\label{EqTriangleInequality}
h(s) \ge  c'' h_{\overline{\mathbb{A}_g},\cM}(\tau(s)) - c'''
\end{equation}
for all $s \in \mathbb{M}_g(\IQbar)$; see
\cite[Lemma~4]{Silverman:heightest11}, the left-hand side is the Weil
height and represents a height function coming from an ample line
bundle on $\overline{\mathbb{M}_g}$.
If $h_{\overline{\mathbb{A}_g},\cM}(\tau(s))\ge c_1$ with $c_1\ge
2 c'''/c''$
then $h(s) \ge c''
h_{\overline{\mathbb{A}_g},\cM}(\tau(s))/2$. We find that it suffices to
prove the alternative with $h_{\overline{\mathbb{A}_g},\cM}(\tau(s))$
replaced by $h(s)$ in (ii) and adjusting $c_3$. 
Our proof is by induction on $\dim S$.

If $\dim S = 0$, then the proposition follows by enlarging $c_1$.

If $\dim S\ge 1$, we fix $M = 3g-2$.
% Let $S^{\mathrm{sm}}$ denote the regular locus of $S$.
Applying
Theorem~\ref{ThmNonDegForBd} to the immersion
 $S^{\mathrm{sm}} \hookrightarrow \mathbb{M}_g$,
%$S \hookrightarrow \mathbb{M}_g$,
we conclude that the closed irreducible subvariety
 $X:=\mathscr{D}_M(\mathfrak{C}_{S^{\mathrm{sm}}}^{[M+1]})$
%$X=\mathscr{D}_M(\mathfrak{C}_{S}^{[M+1]})$
of the abelian scheme
% $\mathfrak{A}_g^{[M]}\times_{\mathbb{A}_g} S^{\mathrm{sm}} \rightarrow
% S^{\mathrm{sm}}$
$\cA=\mathfrak{A}_g^{[M]}\times_{\mathbb{A}_g} S^{\mathrm{sm}} \rightarrow
S^{\mathrm{sm}}$
is non-degenerate. Hence we can apply Theorem~\ref{ThmHtInequality} to
%$\cA= \mathfrak{A}_g^{[M]}\times_{\mathbb{A}_g} S^{\mathrm{sm}}$
$\cA$ %= \mathfrak{A}_g^{[M]}\times_{\mathbb{A}_g} S$
and 
% $X=\mathscr{D}_M(\mathfrak{C}_{S^{\mathrm{sm}}}^{[M+1]})$
$X$
(and the compactification $\overline{S}$ is the Zariski closure of $S$ in $\overline{\mathbb{M}_g}$). 
So, combined with \eqref{EqTriangleInequality}, there exist constants
$c > 0$ and $c'$ as well as a Zariski open dense subset $U$ of
%$X=\mathscr{D}_M(\mathfrak{C}_{S^{\mathrm{sm}}}^{[M+1]})$,
$X$,
satisfying the following property.
%For all $s \in S^{\mathrm{sm}}(\IQbar)$ and
For all $s \in S(\IQbar)$ and
all $P, Q_1,\ldots,Q_M \in \mathfrak{C}_s(\IQbar)$, we have
\begin{equation}\label{EqHtInequalityBdRatPt} 
%c h_{\overline{\mathbb{A}_g},\cM}(\tau(s)) \le
c h(s)\le 
\hat{h}(Q_1-P) + \cdots + \hat{h}(Q_M - P) + c' \quad
\text{if}\quad (Q_1-P,\ldots, Q_M-P) \in U(\IQbar).
\end{equation}

Observe that $\pi(X)=S^{\mathrm{sm}}$, where
%Observe that $\pi(X)=S$, where
%$\pi\colon \mathfrak{A}_{g}^{[M]}\times_{\mathbb{A}_g} S\rightarrow
%S$
$\pi\colon \mathcal{A} \rightarrow S^{\mathrm{sm}}$ 
is the structure
morphism.
Therefore, $S \setminus \pi(U)$ is not Zariski dense in
$S$. %\Pcomment{I don't see why $\pi(U)$ is open, but it doesn't matter.} 
Let $S_1,\ldots, S_r$ be the irreducible components of the
Zariski closure of $S\setminus \pi(U)$ in $S$. Then $\dim S_j \le \dim
S - 1$ for all $j$.

By the induction hypothesis, this proposition holds for all $S_j$. Thus it
remains to prove the conclusion of this proposition for curves above
\begin{equation}
\label{eq:soutsideSi}
s \in S(\IQbar) \setminus \bigcup_{j=1}^r S_j(\IQbar) \subset \pi(U(\IQbar)).
\end{equation}
 First we construct $\Xi_s$ and then
we will show that we are in one of the two alternatives. 

It is convenient to fix a base point $P_s\in
\mathfrak{C}_s(\IQbar)$ and consider $(\mathfrak{C}_s - P_s)^M$
as a subvariety of $\cA_s = \pi^{-1}(s)$. 
% as a curve
%inside $(\mathfrak{A}_g)_{\tau(s)}$ via the Abel--Jacobi map.

Let us set $W = X\setminus U$, it is a Zariski closed and proper
subset of $X$.  By (\ref{eq:soutsideSi}) we find  $W_s\subsetneq X_s = \mathscr{D}_M(\mathfrak{C}_s^{[M+1]})$. 

Let $Z$ be an irreducible component of
$W_s$. 
%Let $j_s \colon \mathfrak{C}_s \rightarrow (\mathfrak{A}_g)_s$ be any Abel--Jacobi embedding.
Consider the set
%\Gcomment{Reworded this paragraph (and the next one) to apply the moduli-space free Lemma on packages.}
\[
\Xi_Z := \{ P \in \mathfrak{C}_s(\IQbar) : ( \mathfrak{C}_s- P)^M
\subseteq Z \}. % \subsetneq \mathscr{D}_m(\mathfrak{C}_s^{m+1}) \}.
\]
%which is well-defined and independent of $P_s$.
Apply Lemma~\ref{LemmaPacketsAlternative3} to $A =
(\mathfrak{A}_g)_{\tau(s)}, C=\mathfrak{C}_s-P_s\subset A$, and $Z$. As $Z\subsetneq
\mathscr{D}_M(\mathfrak{C}_s^{[M+1]})$  we have $\#\Xi_Z \le 84(g-1)$.

Let $\Xi_s = \bigcup_Z \Xi_Z$ where $Z$ runs over all irreducible
components of $W_s$. The number of irreducible components  is
bounded from above in an algebraic family.
So the number of irreducible components of $W_s$
is bounded from above by a number that  is independent of $s$; but it may depend on $W$.  We take
$c_2$ to be such a number multiplied with $84(g-1)$. Thus
$\#\Xi_s \le c_2$ if (\ref{eq:soutsideSi}) and with $c_2$ independent
of $s$ and $P$.

Say $P\in \mathfrak{C}_s(\IQbar)$ and $P\not\in \Xi_s$.
So we are not in case (i) of the proposition.
Then $(\mathfrak{C}_s-P)^M
\not\subseteq W_s$. We want to apply Lemma~\ref{LemmaNAlon} to
$\mathfrak{C}_s-P$ and $W_s$.

%Recall  that $V_1,\ldots,V_t$ are Zariski open affine subsets that cover $\mathbb{A}_g$.
% The image of $s\in \mathbb{M}_g(\IQbar)$ under the Torelli morphism is  
% $\tau(s)\in\mathbb{A}_g(\IQbar)$. 
Recall that the abelian scheme $\mathfrak{A}_{g}$ is embedded in
$\IP_{\IQbar}^n\times \mathbb{A}_g$  over
$\mathbb{A}_g$, \textit{cf.}
$\mathsection$\ref{subsec:univfamily}.
So $\cA$ is embedded in $(\IP_{\IQbar}^n)^M\times S^{\mathrm{sm}}$ over $S^{\mathrm{sm}}$. 
We may identify $\mathfrak{C}_s-P$ with
 a smooth curve in $\IP_{\IQbar}^n$. 
The degree of $\mathfrak{C}_s-P$ as a subvariety of $\IP_\IQbar^n$ 
is bounded independently of $s$ by
Lemma~\ref{lem:uniformdegreebound}(i); applying the Torelli
morphism $\tau$ does not affect the degree. 
% The degree of $\mathfrak{C}_s-P$ as a subvariety of 
% $\IP_\IQbar^n$ equals the degree of $\mathfrak{C}_s$ and hence is bounded from above independently 
% of $s$.
%\Gcomment{Extra comment: We cannot apply
 % Lemma~\ref{lem:uniformdegreebound}, which claims the independency
 % for a particular $P_s$. But this sentence is easy to prove (B\'ezout
 % and/or flatness).} \Pcomment{Yes, we don't need
 % Lemma~\ref{lem:uniformdegreebound} and just need that the degree is
 % uniformly bounded from above in an algebraic family.}
%\Pinlinecomment{New.}
Moreover, $W_s$ is Zariski closed in
$X_s\subset \cA_{s}$. Still holding $s$ fixed we may take $W_s$ as a Zariski
closed subset of $(\IP_{\IQbar}^n)^M$. Being the fiber above $s$ of a
subvariety of $(\IP_{\IQbar}^n)^M\times S^{\mathrm{sm}}$, we find that the degree
of $W_s$ is bounded from above independently of $s$. 
From
Lemma~\ref{LemmaNAlon} we
thus obtain a number $c_4$, depending only  on these bounds and with the
following property. Any subset $\Sigma \subset
\mathfrak{C}_s(\IQbar)$ with cardinality $\ge c_4$ satisfies $(\Sigma-P)^M
\not\subset W_s$. 
It is crucial that $c_4$ is independent of $s$.

We work with $\Sigma = \{Q \in \mathfrak{C}_s(\IQbar) : \hat{h}(Q-P)
\le  h(s)/c_3 \}$ with $c_3 = 2M/c$. 
If $\#\Sigma < c_4$, then we are in alternative (ii) of the proposition. 

Finally, let us assume $\#\Sigma \ge c_4$. % we will show $h_{\overline{\mathbb{A}_g},\cM}(\tau(s))
% < c_1$ with $c_1=3c'/c$. Indeed, 
The discussion above implies
that there exist $Q_1,\ldots,Q_M \in
\Sigma$ such that $(Q_1 - P, \ldots, Q_M-P) \not\in
W_s(\IQbar)$, \textit{i.e.}, $(Q_1 - P, \ldots, Q_M-P)\in U(\IQbar)$.
Thus we can apply \eqref{EqHtInequalityBdRatPt} and obtain
\begin{equation*}
%  \label{eq:hScMbound}
  h(s) \le \frac{1}{c }\left( M\frac{c h(s)}{2M}
    +c' \right) = \frac{1}{2} h(s) + \frac{c'}{c}.
\end{equation*}
Hence
$h(s)\le  2c'/c$. Now (\ref{EqTriangleInequality}) implies
$h_{\overline{\mathbb{A}_g},\cM}(\tau(s)) < c_1$ if $c_1 >
(2c'/c+c''')/c''$.
So this case cannot occur if
$h_{\overline{\mathbb{A}_g},\cM}(\tau(s)) \ge c_1$ and $c_1$ is
sufficiently large. 
\end{proof}

%%% Local Variables:
%%% TeX-master: "main"
%%% End:
